% part: intuitionistic-logic
% chapter: propositions-as-types
% section: reduction

\documentclass[../../../include/open-logic-section]{subfiles}

\begin{document}

\olfileid{int}{pty}{red}

\olsection{الاختزال}

في~!!{derivation}s الاستنباط الطبيعي تكون قاعدة~!!a{introduction} تتبعها
قاعدة~!!a{elimination} التفافًا يمكن ``اختزاله''. فمثلًا، يلغي
استدلال~$\Intro{\land}$ يتبعه~$\Elim{\land}$ أحدهما الآخر، لأن نتيجة
$\Elim{\land}$ تكون دائمًا أحد حدي الاقتران الذي ينتجه~$\Intro{\land}$،
ومقدمتا~$\Intro{\land}$ هما حدا الاقتران نفسيهما. ولذلك، إذا احتجنا إلى
!!a{derivation} لأي من الحدين، أمكننا استعمال ذلك~!!{derivation} مباشرةً
من دون إدخال الاقتران ثم حذفه. ولدينا:
\begin{gather*}
  \bottomAlignProof
  \AxiomC{}
  \RightLabel{$\delta_1$}
  \DeduceC{$!A_1$}
  \AxiomC{}
  \RightLabel{$\delta_2$}
  \DeduceC{$!A_2$}
  \RightLabel{$\Intro{\land}$}
  \BinaryInfC{$!A_1 \land !A_2$}
  \RightLabel{$\Elim{\land}_i$}
  \UnaryInfC{$!A_i$}
  \DisplayProof\bottomAlignProof
  \quad
  \redone
  \quad
  \AxiomC{}
  \RightLabel{$\delta_i$}
  \DeduceC{$!A_i$}
  \DisplayProof
\end{gather*}

يمكننا أن ننظر في علاقة اختزال مماثلة على حدود البرهان المقابلة يقترحها هذا
التبسيط. فإذا كان~$N_1$ حد برهان~$\delta_1$ و$N_2$ حد برهان~$\delta_2$،
قلنا إن حد برهان البرهان الواقع على اليسار يختزل في خطوة واحدة إلى حد برهان
البرهان الواقع على اليمين:
\[
  \proj{i}{\pair{N_1}{N_2}} \redone N_i
\]
وهذه في حساب لامبدا المنمّط قاعدة \emph{اختزال بيتا} للنمط الضربي.

يسمى في الاستنباط الطبيعي زوج من الاستدلالات، مثل~$\Intro{\land}$ يتبعه
$\Elim{\land}$، أي زوج يخضع للاختزال، \emph{قطعًا}. وفي حساب لامبدا المنمّط
يسمى الحد الواقع على يسار~$\redone$ \emph{ردكسًا}، والحد الواقع على يمينه
\emph{مختزَلًا}. وعلى خلاف حساب لامبدا غير المنمّط، حيث لا يعد ردكسًا إلا
$(\lambd[x][N])Q$، يوسَّع نحو حساب لامبدا المنمّط ليشمل الحدود التي تحتوي
$\pair{N}{M}$ و$\proj{i}{N}$ و$\inj[!A]{i}{N}$ و
$\dcase{N}{x_1}{M_1}{x_2}{M_2}$ و$\abort{!A}{N}$، مع الردكسات المقابلة.

ولدينا على نحو مماثل اختزال للفصل:
\begin{gather*}
  \bottomAlignProof
  \AxiomC{}
  \RightLabel{$\delta$}
  \DeduceC{$!A_i$}
  \RightLabel{$\Intro{\lor}$}
  \UnaryInfC{$!A_1 \lor !A_2$}
  \AxiomC{$\Discharge{!A_1}{x_1}$}
  \RightLabel{$\delta_1$}
  \DeduceC{$!C$}
  \AxiomC{$\Discharge{!A_2}{x_2}$}
  \RightLabel{$\delta_2$}
  \DeduceC{$!C$}
  \DischargeRule{$\Elim{\lor}$}{x_1,x_2}
  \TrinaryInfC{$!C$}
  \DisplayProof\bottomAlignProof
  \quad
  \redone
  \quad
  \AxiomC{}
  \RightLabel{$\delta$}
  \DeduceC{$!A_i$}
  \RightLabel{$\delta_i$}
  \DeduceC{$!C$}
  \DisplayProof
\end{gather*}
ويقابل ذلك الاختزال الآتي على حدود البرهان:
\begin{gather*}
  \dcase{\inj[!A_{3-i}]{i}{M}}{x_1}{N_1}{x_2}{N_2}
  \redone \Subst{N_i}{M}{x_i}.
\end{gather*}
حيث~$M$ حد برهان~$\delta$، و$N_i$ حد برهان~$\delta_i$. وهذه قاعدة اختزال
بيتا للأنماط الجمعية. والترميز~$!A_{3-i}$ في الحقن هو حد الفصل المحذوف من
المقدمة. ويعني~$\Subst{N_i}{M}{x_i}$ إحلال~$M$ محل كل وقوع حر للمتغير
$x_i$ في~$N_i$.

ويقابل اختزال النمط الدالي حذف التفاف من~$\Intro\lif$ يتبعه~$\Elim\lif$.
\begin{gather*}
  \bottomAlignProof
  \AxiomC{$\Discharge{!A}{x}$}
  \RightLabel{$\delta$}
  \DeduceC{$!B$}
  \DischargeRule{$\Intro{\lif}$}{x}
  \UnaryInfC{$!A \lif !B$}
  \AxiomC{}
  \RightLabel{$\delta'$}
  \DeduceC{$!A$}
  \RightLabel{$\Elim{\lif}$}
  \BinaryInfC{$!B$}
  \DisplayProof\bottomAlignProof
  \quad
  \redone
  \quad
  \AxiomC{}
  \RightLabel{$\delta'$}
  \DeduceC{$!A$}
  \RightLabel{$\delta$}
  \DeduceC{$!B$}
  \DisplayProof
\end{gather*}
أما في حدود البرهان فليس ذلك سوى اختزال بيتا المعتاد:
\begin{gather*}
  (\lambd[\typeof{x}{!A}][N]M)
  \redone \Subst{N}{M}{x}.
\end{gather*}

وإلى جانب تحويلات الاختزال على~!!{proof}s، يمكننا النظر في تحويلات
التبديل. وهي تنطبق على~!!{proof}s أو براهينها الجزئية التي تنتهي بقاعدة
$\Elim{\lor}$ تتبعها قاعدة~!!{elimination} أخرى، مثلًا:
\begin{multline*}
  \bottomAlignProof
  \AxiomC{}
  \RightLabel{$\delta$}
  \DeduceC{$!A_i$}
  \RightLabel{$\Intro{\lor}$}
  \UnaryInfC{$!A_1 \lor !A_2$}
  \AxiomC{$\Discharge{!A_1}{x_1}$}
  \RightLabel{$\delta_1$}
  \DeduceC{$!C_1 \land !C_2$}
  \AxiomC{$\Discharge{!A_2}{x_2}$}
  \RightLabel{$\delta_2$}
  \DeduceC{$!C_1 \land !C_2$}
  \DischargeRule{$\Elim{\lor}$}{x_1,x_2}
  \TrinaryInfC{$!C_1 \land !C_2$}
  \RightLabel{\Elim{\land}[i]}
  \UnaryInfC{$!C_i$}
  \DisplayProof\bottomAlignProof
  \quad
  \redone
  \\
  \AxiomC{}
  \RightLabel{$\delta$}
  \DeduceC{$!A_i$}
  \RightLabel{$\Intro{\lor}$}
  \UnaryInfC{$!A_1 \lor !A_2$}
  \AxiomC{$\Discharge{!A_1}{x_1}$}
  \RightLabel{$\delta_1$}
  \DeduceC{$!C_1 \land !C_2$}
  \RightLabel{\Elim{\land}[i]}
  \UnaryInfC{$!C_i$}
  \AxiomC{$\Discharge{!A_2}{x_2}$}
  \RightLabel{$\delta_2$}
  \DeduceC{$!C_1 \land !C_2$}
  \RightLabel{\Elim{\land}[i]}
  \UnaryInfC{$!C_i$}
  \DischargeRule{$\Elim{\lor}$}{x_1,x_2}
  \TrinaryInfC{$!C_i$}
  \DisplayProof
\end{multline*}
إذا كانت حدود برهان~$\delta$ و$\delta_1$ و$\delta_2$ هي، على الترتيب،
$M$ و$N_1$ و$N_2$، كانت قاعدة الاختزال المقابلة لحدود البرهان، أي حدود
لامبدا، هي:
\[
  \proj{i}{\dcase{M}{x_1}{N_1}{x_2}{N_2}} \redone
  \dcase{M}{x_1}{\proj{i}{N_1}}{x_2}{\proj{i}{N_2}}.
\]

وبالطبع لا نحذف الالتفافات عند نهايات~!!{proof}s وحدها، بل نحذفها أيضًا
داخل~!!{proof}s بتطبيق تحويل الاختزال على~!!{proof} جزئي ينتهي بالتفاف.
وكذلك ينطبق اختزال بيتا على الحدود الجزئية من حدود لامبدا. فإذا كان~$N$
حدًا جزئيًا من~$M$، وكان~$N\redone N'$، أمكننا استبدال~$N'$ بالحد الجزئي
$N$ في~$M$ لنحصل على~$M'$. ونكتب في هذه الحالة أيضًا~$M\redone M'$.
وإذا وجدت متتالية
$M=M_1\redone M_2\redone M_3\redone\dots\redone M_k=M'$، بما في ذلك
حالة~$k=1$، أي~$M=M'$، كتبنا~$M\red M'$.

يسمى~!!{proof} الذي لا يمكن تطبيق أي تحويل على أي~!!{proof} جزئي منه
\emph{عاديًا}. وعلى نحو مماثل، يسمى حد لامبدا الذي لا يمكن تطبيق أي تحويل
عليه أو على أي حد جزئي منه \emph{عاديًا}. ويسمى حد لامبدا العادي أيضًا
\emph{صيغة عادية}.

\end{document}
